%%%%%%%%%%%%%%%%%%%%%%%%%%%%%%%%%%
% main.tex
%
% Hier bitte nur \title und \author ausfüllen
%
%

\documentclass[a4paper]{scrartcl}

\usepackage[utf8]{inputenc}
\usepackage[german]{babel}
\usepackage{amsfonts}
\usepackage{graphicx}
\usepackage{makeidx}
\usepackage{longtable}
\usepackage{setspace}
\usepackage{listings}
\usepackage{xcolor}
%\onehalfspacing  % anderthalbzeilig



% see http://merkel.zoneo.net/Latex/natbib.php
\usepackage{natbib}
% bibpunct: set up for the reference appearance
% 1. the opening bracket symbol, default = (
% 2. the closing bracket symbol, default = )
% 3. the punctuation between multiple citations, default = ;
% 4. the letter `n' for numerical style, or `s' for numerical superscript
%    style, any other letter for author-year, default = author-year;
% 5. the punctuation that comes between the author names and the year
% 6. the punctuation that comes between years or numbers when common
%    author lists are suppressed (default = ,);
%\bibpunct{[}{]}{,}{n}{}{,}
\bibpunct{[}{]}{,}{}{}{,}

\usepackage{ae,aecompl}
\usepackage[pdfborder={0 0 0}]{hyperref}

%\titlehead{ \\
%Carl von Ossietzky Universität Oldenburg \\
% \\
%}
\titlehead{
	\centering
	\includegraphics[width=0.5\columnwidth]{images/cau-sw} \\
	Technische Fakultät \\
	Institut für Informatik \\
	Arbeitsgruppe Software Engineering
}

\title{Kieker Roadmap}
\author{Reiner Jung}
\date{\today}
\subject{v0.0}

%%%%%%%%%%%%%%%%%%%%%%%%%%%%%%%%%%%
% macors.tex
%
% author: Reiner Jung
%

%% einstellung f�r mehrere grafiken auf einer seite
%
%
%
\renewcommand{\textfraction}{0.05} 
\renewcommand{\topfraction}{0.95} 
\renewcommand{\bottomfraction}{0.95} 
\setcounter{topnumber}{4} 
\setcounter{totalnumber}{5} 
 



%% setzt pfad f�r bilder
%
% Parameter:
%   Neuer Pfad f�r Bilder
%
% initialisierung
\newcommand{\img}{images}
% 
\newcommand{\setimage}[1]{\renewcommand{\img}{#1/images}}

%% setzt pfad f�r subdokumente
%
% Parameter:
%   Neuer Pfad f�r Bilder
%
% initialisierung
\newcommand{\docdir}{.}
% 
\newcommand{\setdocdir}[1]{\renewcommand{\docdir}{#1}}

%% input mit docdir usage
%
% Parameter:
%   Datei die eingebunden werden soll
\newcommand{\use}[1]{\input{\docdir/#1}}

%% \includepicture
% Erlaubt es ein Bild in ein Dokument einzuf�gen
%
% Parameter:
%   1. Werte f�r \includegraphics z.B. die Breite width=11.7cm
%   2. Dateiname der eps-Datei im Unterverzeichnis graphics
%   3. Bildunterschrift
%   4. Label/Referenzname

\newcommand{\incpic}[4]{
\begin{figure}[ht]
\begin{center}
\fbox{
\includegraphics[#1]{\img/#2}}
\caption{#3} \label{#4}
\end{center}
\end{figure}}

%% \includepicture
% Erlaubt es ein Bildschirmabbild in ein Dokument einzuf�gen
% Erzeugt keinen Rahmen
%
% Parameter:
%   1. Werte f�r \includegraphics z.B. die Breite width=11.7cm
%   2. Dateiname der eps-Datei im Unterverzeichnis graphics
%   3. Bildunterschrift
%   4. Label/Referenzname

\newcommand{\incscreen}[4]{
\begin{figure}[hpt]
\begin{center}
\includegraphics[#1]{\img/#2}
\caption{#3} \label{#4}
\end{center}
\end{figure}}


%% \figref
% reference a figure
\newcommand{\figref}[1]{(see fig. \ref{#1})}

%% \secref
% reference a section
\newcommand{\secref}[1]{(see section \ref{#1})}


%% \note
% Erzeugt eine Notiz im Text, welche sp�ter noch entfernt werden muss.
% (von Michael um Autor erweitert, da das auch ganz interessant sein k�nnte...) 
%
% Parameter:
%   1. Autor der Notiz
%   2. Die Notiz

\newcommand{\note}[1]{\paragraph{Note:} \textit{#1} \vskip 2em}

%% action
\newcommand{\actionpre}{}
\newcommand{\actionpost}{}
\newcommand{\actionrole}{}

\newcounter{actioncount}
\setcounter{actioncount}{1}

\newenvironment{action}[1]{
\renewcommand{\actionpre}{}
\renewcommand{\actionpost}{}
\renewcommand{\actionrole}{}
\vspace{3mm}
\noindent
\vline\hspace{2mm}
\begin{minipage}{0.95\textwidth}
\paragraph{Action \arabic{actioncount}: #1}
}{
\begin{description}
\ifthenelse{\equal{\actionrole}{}}{}{\item[Roles:] \actionrole}
\ifthenelse{\equal{\actionpre}{}}{}{\item[Precondition:] \actionpre}
\ifthenelse{\equal{\actionpost}{}}{}{\item[Postcondition:] \actionpost}
\end{description}\end{minipage}\\
\vspace{3mm}\stepcounter{actioncount}}

%% \pre
\newcommand{\pre}[1]{\renewcommand{\actionpre}{#1}}

%% \post
\newcommand{\post}[1]{\renewcommand{\actionpost}{#1}}

%% \role
\newcommand{\role}[1]{\renewcommand{\actionrole}{#1}}

%%
\newcounter{reqcount}
\setcounter{reqcount}{1}
\newcommand{\req}[2]{\paragraph{Requirement \arabic{reqcount}: #1} #2\stepcounter{reqcount}}

%% OWL
\lstdefinelanguage{OWL}
{keywords={owl:AllDifferent,
allValuesFrom,
AnnotationProperty,
backwardCompatibleWith,
cardinality,
Class,
complementOf,
DataRange,
DatatypeProperty,
DeprecatedClass,
DeprecatedProperty,
differentFrom,
disjointWith,
distinctMembers,
equivalentClass,
equivalentProperty,FunctionalProperty,hasValue,imports,incompatibleWith,
intersectionOf,
InverseFunctionalProperty,
inverseOf,
maxCardinality,
minCardinality,
Nothing,
ObjectProperty,
oneOf,
onProperty,
Ontology,
OntologyProperty,
priorVersion,
Restriction,
sameAs,
someValuesFrom,
SymmetricProperty,
Thing,
TransitiveProperty,
unionOf,
versionInfo,
about,
ID,
parseType,
resource,
type,
subClassOf,
owl,rdf,rdfs}, keywordsprefix={owl:,rdf:,rdfs:}, morecomment=[s]{<!--}{-->}, sensitive=true,
}

\lstset{tabsize=4, basicstyle=\sffamily\small, keywordstyle=\bfseries, columns=[r]fullflexible}

\definecolor{keywordA}{RGB}{127,24,156}
\definecolor{keywordB}{RGB}{50,50,255}
\definecolor{keywordC}{RGB}{255,50,50}

\lstdefinelanguage{XTEND}{
alsoletter={:->},
morekeywords=[1]{create,::,:,if,else,then},
keywords=[2]{->},
keywords=[3]{this,null},
keywords=[4]{Integer,List},
keywordstyle=[1]\color{keywordA},
keywordstyle=[2]\color{keywordB},
keywordstyle=[3]\color{keywordB},
keywordstyle=[4]\color{keywordC},
stringstyle=\bfseries,
string=[b]{\"},
}

%% \partner
%
% fields can be empty
%
% 1 = name
% 2 = strasse
% 3 = plz + ort
% 4 = tel
% 5 = fax
% 6 = mail
% 7 = www
\newcommand{\partner}[7]{
\pagebreak{}

\section*{Kontakt}

\vbox{
\hbox{MENGES-Projektpartner\hfill}
\vskip2mm
\hbox{\textbf{#1}\hfill}
\vskip2mm
\hbox{#2\hfill}
\hbox{#3\hfill}
\vskip1mm
\hbox{Tel.: #4}
\hbox{Fax: #5}
\vskip1mm
\hbox{E-Mail: \href{mailto:#6}{\texttt{#6}}}
\hbox{Web: \url{#7}}
}
}


%\hyphenation{}

\makeindex

\begin{document}

\maketitle

\tableofcontents

%%%%%%%%%%%%%%%%%%%%%%%%%%%%%%%%%%%%%%%%%%
%%
%% add content here


%%
\section{Ziele}

\begin{itemize}
\item Instrumentierung für verschiedene Programmiersprachen via Source-Code, Code-Injection, Intermediate-Code
\item Sammlung von Auswertungsverfahren
\item Sammlung von Visualisierungen
\item Modellgetriebene Instrumentierung
\item Klare Schnittstellen und Methodiken für alle Komponenten von Kieker
\end{itemize}

%%
\section{Bestehende Konzepte}

%
\subsection{Kieker 1.2}

Kieker-Basis-Projekt mit der klassischen Kieker Instrumentierungstechnik für Java und der Probe, Storage, Auswertungs-Architektur.

%
\subsection{Monitoring Tool}

Jens hat im Rahmen seiner Dis. ein grafisches Werkzeug gebaut, welches ebenfalls zum Kieker-Kontext gehört.

\begin{itemize}
\item Graphischer Editor für die Kombination von Filtern/Analyseplugins
\item Visualisierung der Ergebnisse
\end{itemize}

%
\subsection{Kanalyse}

\begin{itemize}
\item Pipe-and-Filter-Sprache für die Kombination von Filtern und ggf. die Komposition von Filtern.
\item Ecore-Modell für Plugins
\end{itemize}

%%
\subsection{Projekt-Infrastruktur}

Um ein die einzelnen Projekte zusammenführen zu können und ein einheitliches Framework zu bekommen sind verschiedene Schritte notwendig. Für den Anfang schlage ich deshalb ein gemeinsames Repository für Kieker 2.0 vor in das wir dann die einzelnen Werkzeuge integrieren. Da wir alle hip und modern sind empfiehlt sich hier git zu benutzen.

Im Moment wird zwar Sourceforge benutzt, aber ein integriertes Ticket, Versions und Wiki System können wir auch mit Trac und Git realisieren. Ebenso können wir Hudson hier einbinden. Im Grunde entspricht dieses Setup auch dem Setup der RTSYS Gruppe von der wir dann auch entsprechend Technologietransfer haben können. Und ich glaube dieses Setup entspricht auch der generellen Strategie in der Abteilung.

%%
\section{Mögliche und Geplante Arbeiten}

Auf inhaltlicher Ebene sollten wir uns die einzelen Pakete ansehen und prüfen was davon weiter zu verwenden ist und vor allem wie es in ein Gesamtprojekt eingebracht werden kann ohne gleich alles neu schreiben zu müssen.

Nach einem solchen Bootstrapping, können wir uns dann um Detail kümmern, welche auch als Arbeiten für Studierende ausgelegt werden können. Im Einzelnen sind das:

\begin{itemize}
\item Generisches Analyse-Framework (Design) auf Basis von Jens' und Andrés Arbeiten
\item Metamodell für das Logging-Format sowie Transformationen, die aus Logging-Format-Modellen passende Serialisierer (Reader/Writer) erzeugen \fbox{BA}
\item Implementierung des Frameworks \fbox{mehrere Arbeiten}
\item Meta-Modell/Sprache für das Analyse-Framework \fbox{MA}
\item Graphischen Editor für das Analyse-Framework \fbox{BA/MA}
\item Instrumentierungsbeschreibung und Generator für Kieker-artigen Instrumentierungscode für Java \fbox{BA/MA}
\item Instrumentierungsbeschreibung und Generator für Kieker-artigen Instrumentierungscode für C\# \fbox{BA/MA}
\item Instrumentierungsbeschreibung und Generator für Kieker-artigen Instrumentierungscode für ST/FUP \fbox{BA/MA}
\end{itemize}



% einkommentieren, wenn eine Bibliographie existiert
%\bibliographystyle{natdin}
%\bibliography{references}

% einkommentieren, wenn ein Index gewünscht ist 
\printindex

\end{document}
